\chapter{Cifrado Homomórfico sobre Redes Neuronales}

Las redes neuronales artificiales (RNA) han impulsando avances sin precedentes en campos tan diversos como la visión por computadora, el procesamiento del lenguaje natural, la medicina y las finanzas. Su capacidad para aprender patrones complejos directamente de los datos les permite abordar problemas que sin ellas son intratables. Este éxito masivo ha llevado a una proliferación de aplicaciones donde las RNA procesan volúmenes crecientes de información personal y sensible.

Sin embargo, las RNA plantean serios desafíos en materia de privacidad. El procesamiento de datos confidenciales por RNAs, especialmente en entornos de terceros o servicios en la nube, genera problemas sobre la exposición de información sensible. En este contexto el cifrado homomórfico emerge como una solución. Al permitir realizar cómputos directamente sobre datos cifrados, FHE ofrece un mecanismo robusto para preservar la privacidad en las operaciones de las RNA. La aplicación de FHE a las redes neuronales posibilita la ejecución de inferencias sobre datos sensibles manteniendo la confidencialidad.

En este capítulo, nos centraremos en la adaptación de arquitecturas de RNA comunes, concretamente redes neuronales profundas (DNNs) y redes neuronales convolucionales (CNNs), para operar eficientemente bajo esquemas de cifrado homomórfico.

\section{Redes Neuronales Superficiales}

Las RNA son modelos computacionales cuyo propósito principal es aprender patrones y relaciones complejas a partir de datos, permitiéndoles realizar tareas como clasificación, regresión etc. Dentro de este vasto campo, las redes neuronales superficiales representan una categoría caracterizada por tener una capa oculta entre su capa de entrada y su capa de salida.

\subsection{Definición Formal y Estructura General}\label{sec:shallow_general_case}
Hemos descrito redes superficiales de ejemplo para desarrollar la intuición sobre su funcionamiento. Ahora, definimos una ecuación general para una red neuronal superficial $y = f[x, \phi]$ que mapea una entrada multi-dimensional $x \in \mathbb{R}^{D_i}$ a una salida multi-dimensional $y \in \mathbb{R}^{D_o}$ utilizando $D$ unidades ocultas.

Cada unidad oculta $h_d$ se calcula como:
\begin{equation}
h_d = a\left[\theta_{d0} + \sum_{i=1}^{D_i} \theta_{di}x_i\right],
\label{eq:hidden_unit}
\end{equation}

y estas se combinan linealmente para crear la salida:
\begin{equation}
y_j = \phi_{j0} + \sum_{d=1}^{D} \phi_{jd}h_d,
\label{eq:output_unit}
\end{equation}

donde $a[\cdot]$ es una función de activación no lineal. El modelo posee parámetros $\phi = \{\theta_{\bullet\bullet}, \phi_{\bullet\bullet}\}$.

Para ilustrar este caso general, consideremos el ejemplo de una red que mapea una entrada escalar $x$ a una salida escalar $y$, y que posee diez parámetros $\phi = \{\phi_0, \phi_1, \phi_2, \phi_3, \theta_{10}, \theta_{11}, \theta_{20}, \theta_{21}, \theta_{30}, \theta_{31}\}$:
\begin{equation}
y = f[x, \phi] = \phi_0 + \phi_1 a[\theta_{10} + \theta_{11}x] + \phi_2 a[\theta_{20} + \theta_{21}x] + \phi_3 a[\theta_{30} + \theta_{31}x].
\label{eq:shallow_nn_example}
\end{equation}

Este cálculo puede desglosarse en tres partes: primero, se computan tres funciones lineales de los datos de entrada ($\theta_{10} + \theta_{11}x$, $\theta_{20} + \theta_{21}x$, y $\theta_{30} + \theta_{31}x$). Segundo, los tres resultados se pasan a través de una función de activación $a[\cdot]$. Finalmente, las tres activaciones resultantes se ponderan con $\phi_1, \phi_2$, y $\phi_3$, se suman, y se añade un sesgo $\phi_0$.

La función de activación permite al modelo describir relaciones no lineales entre la entrada y la salida; como tal, debe ser no lineal. Con una función de activación nula o lineal, el mapeo general de la entrada a la salida estaría restringido a ser lineal. Se han probado muchas funciones de activación diferentes, pero la elección más común es la unidad lineal rectificada o ReLU:
\begin{equation}
a[z] = \text{ReLU}[z] = 
\begin{cases} 
0 & z < 0 \\
z & z \geq 0 
\end{cases}
\label{eq:relu_definition}
\end{equation}

Esta función devuelve la entrada cuando es positiva y cero en caso contrario. Con activaciones ReLU, la red divide el espacio de entrada en politopos convexos definidos por las intersecciones de hiperplanos calculados por las "articulaciones" en las funciones ReLU. Cada politopo convexo contiene una función lineal diferente. Los politopos son los mismos para cada salida, pero las funciones lineales que contienen pueden diferir.

\subsection{Entradas y Salidas Multivariadas}

En el ejemplo anterior, la red tenía una única entrada escalar $x$ y una única salida escalar $y$. Sin embargo, el teorema de aproximación universal también es válido para el caso más general en el que la red mapea entradas multivariadas $x = [x_1, x_2, \dots, x_{D_i}]^T$ a predicciones de salida multivariadas $y = [y_1, y_2, \dots, y_{D_o}]^T$.

Para extender el modelo y predecir salidas multivariadas $y$, simplemente utilizamos una función lineal diferente de las unidades ocultas para cada salida. Por ejemplo, una red con una entrada escalar $x$, cuatro unidades ocultas $h_1, h_2, h_3, h_4$, y una salida multivariada 2D $y = [y_1, y_2]^T$ se definiría como:
\begin{align*}
    h_1 &= a[\theta_{10} + \theta_{11}x] \\
    h_2 &= a[\theta_{20} + \theta_{21}x] \\
    h_3 &= a[\theta_{30} + \theta_{31}x] \\
    h_4 &= a[\theta_{40} + \theta_{41}x],
    \tag{\ref{eq:hidden_units_multivariate}}
\end{align*}
y las salidas serían:
\begin{align*}
    y_1 &= \phi_{10} + \phi_{11}h_1 + \phi_{12}h_2 + \phi_{13}h_3 + \phi_{14}h_4 \\
    y_2 &= \phi_{20} + \phi_{21}h_1 + \phi_{22}h_2 + \phi_{23}h_3 + \phi_{24}h_4.
    \tag{\ref{eq:outputs_multivariate}}
\end{align*}
Las dos salidas son funciones lineales diferentes de las mismas unidades ocultas.

Como se observa, los "puntos de inflexión" en las funciones por partes dependen de dónde las funciones lineales iniciales $\theta_{\bullet0} + \theta_{\bullet1}x$ son "recortadas" por las funciones ReLU $a[\cdot]$ en las unidades ocultas. Dado que ambas salidas $y_1$ e $y_2$ son funciones lineales diferentes de las mismas cuatro unidades ocultas, los cuatro "puntos de inflexión" en cada una deben estar en los mismos lugares. Sin embargo, las pendientes de las regiones lineales y el desplazamiento vertical general pueden diferir.

\subsection{Terminología}
\subsubsection{Neurona o Unidad de Procesamiento}

Una neurona artificial, o unidad de procesamiento, es el componente fundamental de una red neuronal. Recibe una o más entradas, realiza una operación (generalmente una suma ponderada de sus entradas más un sesgo), y luego aplica una función de activación no lineal al resultado para producir una salida.

\subsubsection{Capas: Entrada, Oculta y Salida}
Las neuronas en una red neuronal se organizan en capas.
\begin{itemize}
    \item \textbf{Capa de Entrada:} Recibe los datos iniciales o características del problema.
    \item \textbf{Capa Oculta:} Son capas intermedias entre la entrada y la salida.
    \item \textbf{Capa de Salida:} Produce el resultado final de la red.
\end{itemize}

\subsubsection{Pesos y Sesgos}
\begin{itemize}
    \item \textbf{Pesos ($W$):} Son parámetros ajustables que representan la fuerza de la conexión entre las neuronas de capas adyacentes.
    \item \textbf{Sesgos ($b$):} Son parámetros ajustables que permiten a una neurona activar su función de activación incluso si todas sus entradas son cero.
\end{itemize}

\subsubsection{Función de Activación}
Una función de activación introduce no linealidad en la salida de una neurona. Es crucial para permitir que la red aprenda relaciones complejas y no lineales en los datos. Sin una función de activación no lineal, una red neuronal, independientemente de su profundidad, se comportaría como un modelo lineal. Ejemplos comunes incluyen ReLU, Sigmoide y Tanh.

\section{Redes Neuronales Profundas}
    \subsection{Composición de Redes Neuronales}
    \subsection{De la Composición a las Redes Profundas}
    \subsection{Notación Matricial}

\section{Redes Neuronales Convolucionales}
    \subsection{Introducción a las CNNs y su Motivación}
    \subsection{Capas Convolucionales}
        \subsubsection{Filtros (Kernels), Stride y Padding}
        \subsubsection{Operación de Convolución}
        \subsubsection{Compartición de Pesos y Conectividad Local}
    \subsection{Capas de Agrupamiento (Pooling Layers)}
        \subsubsection{Max Pooling}
        \subsubsection{Average Pooling}
    \subsection{Arquitectura Típica de una CNN}
    \subsection{Ventajas y Aplicaciones de las CNNs}


\section{Operaciones Lineales en CKKS}

La implementación de operaciones lineales, como el producto escalar, la multiplicación matriz-vector y las convoluciones, sobre datos cifrados con el esquema CKKS \cite{CKKS17} presenta desafíos únicos que requieren un diseño cuidadoso. A diferencia de la computación en texto plano, donde las operaciones se realizan directamente sobre los valores numéricos, en el dominio homomórfico debemos trabajar con polinomios cifrados que representan vectores de números complejos. La eficiencia y la viabilidad de estas operaciones dependen en gran medida de cómo se organizan los datos dentro de los textos cifrados y cómo se aprovechan las propiedades del esquema CKKS.

\subsection{El Empaquetado de Datos}

El esquema CKKS permite empaquetar múltiples números en un solo texto cifrado. Esto significa que un único polinomio cifrado puede contener un vector de $N/2$ valores independientes, donde $N$ es el grado del polinomio ciclotómico. La clave para la eficiencia en la inferencia de redes neuronales homomórficas reside en cómo se transforman los tensores (vectores de entrada, matrices de pesos) de la red a las ranuras de los textos cifrados.

El desafío principal radica en que las operaciones homomórficas aumentan el ruido y consumen niveles de la jerarquía de módulos. Por lo tanto, es crucial minimizar el número de operaciones, en particular las rotaciones, que son computacionalmente costosas y requieren claves de evaluación adicionales (Llaves de Galois). Un algoritmo para operar los datos ineficiente puede llevar a un número excesivo de rotaciones, lo que degrada significativamente el rendimiento.

Consideremos una operación fundamental como la multiplicación matriz-vector $y = Wx$. Si $W$ es una matriz de $m \times n$ y $x$ es un vector de $n$ elementos, el resultado $y$ es un vector de $m$ elementos. La forma en que se empaquetan $W$ y $x$ en los textos cifrados determinará la complejidad de la operación. Un enfoque sencillo de empaquetar cada fila o columna de $W$ en un texto cifrado separado resultaría en un gran número de textos cifrados y una gestión compleja. Por ello, se buscan estrategias que permitan empaquetar la mayor cantidad de información posible en un número reducido de textos cifrados, aprovechando al máximo la capacidad SIMD de CKKS.

\subsection{Producto Punto}

El producto punto de dos vectores, $a$ y $b$, se define como la suma de los productos de sus componentes correspondientes: $\sum a_i b_i$. En el contexto del cifrado homomórfico CKKS, esta operación se ejecuta usando rotaciones.

El algoritmo para calcular el producto punto homomórfico de dos vectores cifrados $ct_a$ y $ct_b$ se realiza en dos fases principales:

\begin{enumerate}
    \item Multiplicación Elemento a Elemento:
        Se realiza una multiplicación homomórfica directa entre $ct_a$ y $ct_b$. Dado que CKKS opera sobre vectores empaquetados en ranuras, esta operación produce un nuevo texto cifrado, $ct_{prod}$, donde cada *ranura* contiene el producto $a_i \cdot b_i$ de los elementos correspondientes.

    \item Suma de los Elementos:
        Una vez obtenidos todos los productos $a_i \cdot b_i$ en las ranuras de $ct_{prod}$, el siguiente paso es sumarlos para obtener el resultado final del producto punto. El texto cifrado se rota $ct_{prod}$ a la mitad de sus ranuras y se suma a si mismo, luego se rota el resultado a un cuarto de sus ranuras y se suma, así sucesivamente, hasta que la suma de todos los elementos se acumule en la primera ranura. Estas rotaciones son optimizadas por técnicas como las desarrolladas por Halevi y Shoup \cite{Halevi14}.
\end{enumerate}


\subsection{Multiplicación Matriz Vector}

Para optimizar la multiplicación matriz-vector en el dominio homomórfico, una de las técnicas más eficientes es la diagonalización de la matriz de pesos \cite{Halevi14}, popularizada por Halevi y Shoup. Esta técnica se basa en la observación de que la multiplicación de un vector por una matriz puede expresarse como una suma de productos elemento a elemento \cite{Goodfellow16} con rotaciones cíclicas.

La idea central es transformar la matriz de pesos $W$ en un conjunto de vectores que representan sus diagonales. Específicamente, para una matriz $W$, se extraen sus diagonales y se codifican en textos claros que luego se utilizarán en la operación. Cuando se multiplica un vector cifrado $x$ por una matriz $W$, cada elemento del vector resultante $y_i$ es el producto escalar de la $i$-ésima fila de $W$ con el vector $x$. El proceso seria el siguiente:

\begin{enumerate}
    \item Se codifica el vector de entrada $x$ en un texto cifrado $ct_x$.
    \item La matriz de pesos $W$ se pre-procesa para extraer sus diagonales. Cada diagonal se codifica sin cifrar.
    \item La operación se realiza iterando sobre las diagonales de $W$. Para cada diagonal, se multiplica homomórficamente con una versión rotada de $ct_x$.
    \item Los resultados de estas multiplicaciones y rotaciones se acumulan mediante sumas homomórficas para obtener el texto cifrado del vector resultante $y$.
\end{enumerate}

Un enfoque simple para la multiplicación matriz-vector, sin la diagonalización, implicaría un número significativamente mayor de rotaciones. Para una matriz $W$ de $m \times n$ y un vector $x$ de $n$ elementos, una implementación directa podría requerir $O(m \times n)$ rotaciones para alinear correctamente cada elemento de $W$ con los de $x$ antes de la multiplicación y suma. En contraste, la técnica de diagonalización reduce drásticamente este costo a solo $n$ rotaciones (una por cada diagonal), $n$ multiplicaciones elemento a elemento y $n$ sumas, lo que representa una mejora de rendimiento considerable.

Esta técnica es particularmente eficiente porque reduce el número de rotaciones necesarias. En lugar de rotar el texto cifrado de entrada para cada elemento de la matriz, se rotan los resultados intermedios o se utilizan rotaciones específicas que se corresponden con las diagonales de la matriz. Esto minimiza el consumo de niveles y el costo computacional, haciendo que la multiplicación matriz-vector sea práctica en CKKS.

\subsection{Convoluciones}

Las operaciones de convolución son fundamentales en las Redes Neuronales Convolucionales (CNNs) y, al igual que la multiplicación matriz-vector, requieren una adaptación específica para ser ejecutadas en el dominio homomórfico. La estrategia más común y eficiente para realizar convoluciones homomórficas es transformarlas en una multiplicación matriz-vector, una técnica conocida como `im2col` \cite{Goodfellow16}.

El proceso se descompone en los siguientes pasos:

\begin{enumerate}
    \item Transformación del Tensor de Entrada:
    \begin{itemize}
        \item El tensor de entrada se transforma en una matriz 2D. Cada seccion de la entrada que sería procesado por el kernel de convolución se transforma en una columna de esta nueva matriz.
        \item Esta transformación se realiza en texto plano antes de cifrar los datos, o se puede adaptar para operar sobre datos cifrados si la entrada ya está cifrada.
        \item Los parámetros de la convolución, como el tamaño del kernel, el stride y el padding, se incorporan en esta transformación `im2col`.
    \end{itemize}

    \item Transformación del Kernel de Convolución:
    \begin{itemize}
        \item El kernel (o filtro) de convolución, que es un tensor pequeño, también se transforma en una matriz 2D. Esta matriz se organiza de manera que cada fila corresponda a un kernel aplanado.
    \end{itemize}

    \item Multiplicación Matriz-Vector Homomórfica:
    \begin{itemize}
        \item Una vez que el tensor de entrada y el kernel se han transformado en matrices 2D, la operación de convolución se convierte en una multiplicación matriz-vector estándar.
        \item Se aplica el algoritmo para multiplicación Matriz-Vector descrito en la Sección 1.3. La matriz `im2col` de la entrada se trata como el vector de entrada cifrado, y la matriz del kernel se trata como la matriz de pesos.
        \item El resultado de esta multiplicación es un texto cifrado que representa el mapa de características de salida aplanado.
    \end{itemize}
\end{enumerate}

\section{Funciones de Activación}

Uno de los mayores desafíos al implementar redes neuronales sobre cifrado homomórfico es el manejo de las funciones de activación no lineales. El esquema CKKS solo soporta de forma nativa adiciones y multiplicaciones. Funciones como Sigmoide, Tanh o ReLU no pueden ser evaluadas directamente.

La solución a este problema es aproximar estas funciones de activación mediante polinomios. La elección del polinomio de aproximación es un compromiso crítico entre la precisión del modelo y el rendimiento del cómputo homomórfico.

\subsection{Técnicas de Aproximación}

Existen varias estrategias para encontrar una aproximación polinómica adecuada:

\begin{itemize}
    \item Aproximación por Series de Taylor: Para funciones suaves como la sigmoide o la tangente hiperbólica, se puede utilizar su expansión en series de Taylor alrededor de un punto. Esta es una técnica directa sugerida en los trabajos fundacionales del esquema CKKS \cite{CKKS17}.

    \item Uso de Polinomios Simples: En muchos casos, especialmente para reducir la profundidad del circuito, se reemplaza la función de activación por un polinomio de bajo grado. Una de las aproximaciones más comunes y eficientes es usar la función cuadrática $f(x) = ax^2 + bx + c$. Por ejemplo, la red CryptoNets \cite{GiladBachrach2016CryptoNets} popularizó el uso de la función cuadrada $f(x) = x^2$ como una alternativa a ReLU. Aunque menos precisa, esta aproximación introduce la no linealidad necesaria con un costo computacional mínimo (una sola multiplicación homomórfica).
\end{itemize}

\subsection{Impacto en la Profundidad del Circuito}

Cada multiplicación homomórfica consume un nivel en el esquema CKKS, reduciendo el presupuesto de ruido y el módulo del texto cifrado. La evaluación de un polinomio de grado $d$ requiere una profundidad multiplicativa de $O(\sqrt d)$, utilizando algoritmos de evaluación paralela como el propuesto por Paterson y Stockmeyer \cite{Paterson1973}. Por lo tanto, un polinomio de aproximación de alto grado, aunque más preciso, aumentará significativamente la profundidad del circuito homomórfico. Esto exige parámetros criptográficos más grandes (mayor tamaño de módulo inicial), lo que a su vez incrementa el costo computacional y el tamaño de las claves.

La elección del polinomio de aproximación es, por tanto, una decisión de diseño fundamental en la inferencia de redes neuronales homomórficas, donde se debe balancear la precisión del modelo con las limitaciones y costos del esquema de cifrado \cite{benaissa2021tenseal}.
